\section{Algebraicity with based holonomy markings}
\label{sec:algebraicity}

An algebraic power series will mean an element of
\(\mathbb C\langle x_1,\ldots,x_N\rangle\), the henselization of
\(\mathbb C[x_1,\ldots,x_N]_{(x)}\) inside \(\mathbb C[[x]]\).
We use the edge manifold, boundary coordinates, complex \(C\), and potential
\(f\) of Proposition~\ref{prop:edge-potential}.  All ideals in this section
are scheme-theoretic ideals; no passage to radicals is understood.

\Needspace{12\baselineskip}
\begin{theorem}[Algebraic edge potential with markings]
\label{thm:edge-algebraicity}
Put \(N=\dim(H^1(C)\otimes\mathfrak{sl}_2)\),
\(R=\mathbb C[[h_1,\ldots,h_N]]\), and \(J=\Jac(f)\).
There are an invertible formal coordinate change \(x=\varphi(h)\) and an
algebraic power series \(F(x)\in(x)^3\) such that
\[
                         f(h)=F(\varphi(h)).
\]
Fix a torus basepoint and any finite collection of based loops in \(Y\).
For every matrix entry of their holonomies on the central critical scheme
\(\Spf(R/J)\), there is an algebraic power series in the \(x\)-coordinates
whose pullback is that entry.  In particular this holds for the based
\(\beta\)-holonomy.  The coordinate change can be chosen to identify these
markings on the entire, possibly nonreduced, critical scheme.
\end{theorem}

The proof has five steps.  First, the punctured filling provides an
algebraic Darboux chart.  Second, the full relative closed-form comparison
identifies its symplectic class with the integrated trace on the cochain
model.  Third, transfer gives the Hamiltonian identity for the given
potential \(f\).  Fourth, the closed-form homotopy expresses the
difference of the two potentials as an element of \(\Jac(f)^2\).
Fifth, a formal coordinate correction removes that difference while
fixing the critical scheme and every marked based holonomy.

\subsection{The punctured filling and its algebraic critical chart}

Let \(M\) be a solid torus with an open ball removed.  Glue its torus boundary
to \(\partial Y\), with the meridian glued to \(\alpha\), and write
\[
                   Z=Y\cup_{T^2}M,\qquad \partial Z=S^2.
\]
Choose a point on this sphere and a path to the chosen torus basepoint.
The filling kills \(\alpha\) in the fundamental group.  Thus its based
representations are the based representations of \(\pi_1Y\) with
\(G_\alpha=1\), and the path identifies their torus markings.
The sphere boundary will supply the relative orientation used below.
Set \(G=SL_2\).  For a finite CW complex \(K\), write
\(\Loc_G(K)=\operatorname{Map}(K_B,BG)\); a superscript \(\mathrm{fr}\)
means that the local system is trivialized at the chosen basepoint.
These are derived mapping stacks.  In particular, framing is a derived
fiber of evaluation at the basepoint, not a quotient by conjugation.

The invariant pairing \(\operatorname{tr}(XY)\) is nondegenerate on
\(\mathfrak g=\mathfrak{sl}_2\) and defines a \(2\)-shifted symplectic form
on \(BG\).  Betti transgression gives a \(0\)-shifted symplectic form on
\(\Loc_G(S^2)\) \cite[Theorem~2.5 and Corollary~2.6(4)]{PTVV2013}.
Restriction
\[
                     \Loc_G(Z)\longrightarrow\Loc_G(S^2)
\]
is Lagrangian \cite[Definition~2.6, Claim~2.7, and Theorem~2.9]{Calaque2015}.
Here the hypotheses of these constructions have concrete meanings.
The manifolds have finite CW type, so the Betti sources are
\(\mathcal O\)-compact and their mapping stacks are derived Artin stacks
locally of finite presentation.  Integration over the fundamental class
is the orientation map.  Poincar\'e--Lefschetz duality supplies the required nondegeneracy for
\emph{every} locally constant perfect complex \(E\) over every base
cdga \(D\), including homotopy-coherent coefficient systems:
\[
 C^*(Z,S^2;E)\simeq
 \mathbf R\!\operatorname{Hom}_D\bigl(C^*(Z;E^\vee),D\bigr)[-3].
\]
To see the coefficient scope, local derived duality on contractible
stars identifies the oriented interior dualizing complex with \(D[3]\)
and takes \(E\) to \(E^\vee\otimes D[3]\).  These natural local
identifications descend with all higher restriction coherences;
relative boundary cells give the displayed equivalence.  The finite
constructions in perfect fibers commute with derived base change.
This is the compact-manifold application of
\cite[\S\S3.1.1--3.1.2]{Calaque2015}, and includes adjoint coefficients.
The isotropic homotopy is integration over the relative fundamental chain
of \((Z,S^2)\).

The second Lagrangian is
\[
       \Loc_G(S^2)^{\mathrm{fr}}\longrightarrow\Loc_G(S^2).
\]
Indeed, the based sphere stack is
\(\{1\}\mathbin{\times^{\mathbf R}_G}\{1\}\), the derived
self-intersection of the identity point in the smooth group \(G\).
It is the affine odd vector space with functions
\(\Sym(\mathfrak g^*[1])\), whose generators have degree \(-1\).
On this scheme a form of exterior degree \(2+j\) has internal degree at
most \(-2-j\).  Thus neither internal degree \(-j\), used by a
\(0\)-shifted closed two-form, nor internal degree \(-j-1\), used by its
nullhomotopy, occurs.  The pulled-back form is zero, and its isotropic
path is unique up to homotopy.  At the unique classical point, the tangent
complex of the unframed sphere stack is
\(C^*(S^2;\mathfrak g)[1]\); the framed tangent retains exactly the
degree-two cohomology summand.  Integration and trace pair this summand
perfectly with the omitted degree-zero summand.  This proves the
Lagrangian nondegeneracy condition.  There are no other classical points
at which to check it.

The shifted intersection theorem therefore equips
\begin{equation}
\label{eq:framed-filling-intersection}
 \mathcal X
 =\Loc_G(Z)\mathbin{\times^{\mathbf R}_{\Loc_G(S^2)}}
                   \Loc_G(S^2)^{\mathrm{fr}}
 =\Loc_G(Z)^{\mathrm{fr}}
\end{equation}
with a \((-1)\)-shifted symplectic structure
\cite[Theorem~2.9]{PTVV2013}.  This is a derived scheme.  To see the
finite-presentation assertion directly, a compact connected
three-manifold with nonempty boundary has a finite spine of dimension at
most two: a handle decomposition relative to a collar can be chosen
with no absolute three-handles.  Collapsing a maximal tree of its
one-skeleton gives one vertex, say \(r\) edges, and \(s\) two-cells.
Consequently its framed local-system scheme is the derived fiber
\[
                    G^r\mathbin{\times^{\mathbf R}_{G^s}}\{1\},
\]
where the map records the attaching words.  This description also shows
that every fixed based holonomy is a regular \(G\)-valued function on
the classical scheme.

The algebraic Darboux theorem applies over \(\mathbb C\), at the trivial
classical point of \(\mathcal X\).  More precisely,
\cite[Theorem~5.18(i)]{BBJ2019} supplies a Zariski local algebraic
Darboux chart minimal at that point; in shift \(-1\),
\cite[Example~5.15]{BBJ2019} presents its algebra as
\begin{equation}
\label{eq:algebraic-darboux-chart}
 A^0[p_1,\ldots,p_N],\qquad |p_i|=-1,
 \qquad Qp_i=\frac{\partial F}{\partial x_i},
\end{equation}
where \(A^0\) is smooth and the \(x_i\) are \(\acute{e}\)tale coordinates.
Minimality means that the cotangent differential vanishes at the point
\cite[Definition~2.13]{BBJ2019}; in this presentation it is the Hessian
of \(F\).  Subtracting the constant term therefore gives \(F\in(x)^3\).
In the completed \(x\)-coordinates this is an algebraic power series.
Each of the finitely many regular holonomy entries on the classical
critical chart lifts through the quotient \(A^0\to A^0/\Jac(F)\).
After shrinking the smooth chart, the lifts are regular, and hence give
algebraic power series in the same coordinates.  This assertion concerns
the matrix entries of holonomy; it makes no assertion that their formal
logarithms are algebraic.

The Darboux chart supplies an algebraic potential and regular holonomy
functions on its critical scheme.  To identify that potential with the
particular Chern--Simons function \(f\), we need more than an
isomorphism of critical schemes or equality of tangent pairings.
The full relative closed-form comparison below will give a difference
of potentials in \(\Jac(f)^2\).  A coordinate correction can then remove
that difference while acting trivially on the marked critical scheme.
We state the comparison here and give its proof in
Appendix~\ref{app:formal-trace-comparison}, so that we can first follow
its application to the edge complex.

\subsection{Comparison of the formal trace forms}
\label{subsec:formal-trace-statement}

The following comparison identifies the full closed form on the cochain
model, including the relative integration homotopy.  Write
\(\kappa(u,v)=\operatorname{tr}(uv)\), and denote its \(2\)-shifted
closed form on \(BG\) by \(\omega_\kappa\).
For a nonpositive augmented completed semifree coefficient cdga \((R,Q)\), let
\(\Omega_R^p\) mean continuous coordinate forms, with their internal
grading.  These compute the exterior powers of the cotangent complex.
If \(A_X\) is a commutative cochain model of a finite Betti source, put
\[
 \operatorname{Cl}^2_R(A_X)^n
   =\prod_{j\geq0}
        \bigl(A_X\widehat\otimes\Omega_R^{2+j}\bigr)^{n-j}.
\]
Its differential combines \(d_X+\mathcal L_Q\) and coordinate de Rham
differentiation \(\delta_R\), with the total-complex signs.
The factors of \(A_X\) have exterior weight zero.
Coordinate differentiation is extended by
\(\delta_R(c\otimes u)=(-1)^{|c|}c\otimes\delta_Ru\).
In particular a strict representative has only its \(j=0\) component.
Completion in this notation is in the coefficient augmentation ideal;
the displayed product is the separate completion in exterior weight.

\Needspace{10\baselineskip}
\begin{proposition}[Formal trace comparison]
\label{prop:formal-trace-comparison}
Let \(X\) be a finite triangulated space, and let
\[
 a_X\in
  \bigl(A_X\otimes\mathfrak g\widehat\otimes\mathfrak m_R\bigr)^1,
 \qquad
 (d_X+Q)a_X+\tfrac12[a_X,a_X]=0,
\]
represent a family of local systems formally near the trivial one.
For its evaluation map
\(\operatorname{ev}_{a_X}:X_B\times\Spf R\to BG\), the image of
\(\operatorname{ev}_{a_X}^*\omega_\kappa\) under the Betti
coefficient map of \cite[\S2.1]{PTVV2013} is the class of
\begin{equation}
\label{eq:formal-trace-representative}
 \Theta_X=\tfrac12\kappa(\delta_Ra_X,\delta_Ra_X),
 \qquad
 (d_X+\mathcal L_Q)\Theta_X=0,\qquad \delta_R\Theta_X=0,
\end{equation}
in the full complex \(\operatorname{Cl}^2_R(A_X)\).
This identification is natural for source restrictions, derived
coefficient changes, and integration of source cochains.

For an oriented compact three-manifold \(Z\) with boundary \(S\), it
also identifies the relative-orientation isotropic homotopy with
integration of \(\Theta_Z\) over \(Z\).  If \(S=S^2\) and the sphere
condition is represented by \(a_S=\nu b\), where \(\nu\) is a
two-form and \(b\) has internal degree \(-1\), then
\(\Theta_S=0\) strictly.  Using the zero isotropic path on this sphere
model, the resulting intersection form is represented by
\[
                 \tfrac12\int_Z
                     \kappa(\delta_Ra_Z,\delta_Ra_Z)
\]
as a closed two-form of degree \(-1\).
\end{proposition}

The proof is given in Appendix~\ref{app:formal-trace-comparison}.
It constructs the comparison on the formal group nerve, extends it to
derived coefficient algebras, and checks source evaluation and relative
integration in the full closed-form complex.  Throughout the comparison,
the marked holonomies are the forward transports fixed in
Equation~\eqref{eq:edge-holonomy-sign}.


\subsection{A model preserving the integrated trace}

We identify the cochain model of
\eqref{eq:framed-filling-intersection} with the edge complex \(C\).
On \(S^2\), choose a closed area form \(\nu\).  The inclusion
\(\mathbb C\nu\) into the augmented-kernel forms on the based sphere is
a quasi-isomorphism.  Thus the sphere boundary condition may be imposed
by requiring its trace to lie in \(\mathbb C\nu\), with zero trace in
the other degrees.  Restriction of unrestricted forms to a boundary is
surjective, so this strict fiber computes the required homotopy fiber.

The representatives on the punctured solid torus can be chosen as
follows.  Extend \(\xi\) using a bump density in its longitudinal
circle coordinate.  Put the puncture and a tube from the support disk
of \(\eta\) to the puncture in the complement of that longitudinal
strip.  A transverse Thom form of the tube extends \(\eta\) to a
closed two-form \(\widetilde\eta\) whose sphere trace is the appropriately
oriented \(\nu\).  Denote the extended one-form by
\(\widetilde\xi\).  These choices give
\begin{equation}
\label{eq:zero-added-pairing}
 d\widetilde\xi=d\widetilde\eta=0,
 \qquad \widetilde\xi\wedge\widetilde\eta=0.
\end{equation}
The span of these two forms, with zero multiplication, is a
quasi-isomorphic subalgebra of the forms on \(M\) with the sphere
boundary condition.  In fact that complex has one-dimensional
cohomology in degrees one and two and no other cohomology, as follows
from the relative sequence for \((M,S^2)\); the displayed forms
represent the two generators.

For gluing, use forms smooth on each piece with matching tangential
traces.  This piecewise smooth de Rham model computes the cochains of
\(Z\); its integration is the sum over the two pieces, and their torus
boundary terms cancel by Stokes.  Now glue the small model to the
unrestricted forms on \(Y\).
Surjectivity of restriction from \(Y\) to the torus shows that replacing
the model on \(M\) by this quasi-isomorphic subalgebra preserves the
homotopy pullback.  The resulting algebra is exactly \(C\): its torus
traces are zero in degree zero, multiples of \(\xi\) in degree one,
and multiples of \(\eta\) in degree two.  Moreover,
\eqref{eq:zero-added-pairing} shows that every complementary-degree
pairing on the added piece vanishes before integration.  The full
integrated pairing on the filling is consequently
\begin{equation}
\label{eq:edge-cyclic-pairing}
                \langle u,v\rangle
                  =\int_Y\operatorname{tr}(u\wedge v)
\end{equation}
on \(C\otimes\mathfrak g\).  The comparison of full closed forms and
of the relative integration homotopy in
Proposition~\ref{prop:formal-trace-comparison} applies to this
replacement.  It identifies the shifted symplectic structure of
\(\mathcal X\), rather than just its value on tangent cohomology, with
this strict integrated form.

Choose the basing path through \(M\) outside the longitudinal strip.
The degree-one connection on \(M\) is a multiple of
\(\widetilde\xi\), so its transport along this path is the identity.
The resulting comparison uses the fixed torus framing and preserves
based holonomies as matrices, including the chosen \(\beta\)-holonomy.

\subsection{Transfer with an arbitrary cochain contraction}

The preceding model identifies the integrated trace form and the based
holonomies. We next transfer it to the finite coordinate space used by
the edge slice, keeping the original contraction. The required output
is the Hamiltonian identity for the particular potential $f$.

Write \(H^j=H^j(C)\otimes\mathfrak g\), and retain the splitting
\((i,\operatorname{pr},K)\) from the edge construction, so
\[
 dK+Kd=1-i\operatorname{pr},\qquad
 K^2=Ki=\operatorname{pr}K=0.
\]
No compatibility of \(K\) with the pairing is required.  Choose bases
of \(H^1,H^2\), and let \(h\) and \(q\) be their corresponding
coordinates of degrees zero and \(-1\).  In the completed semifree
algebra \(R[q_1,\ldots,q_N]\), set
\begin{equation}
\label{eq:transfer-v}
 s=i_1h-\tfrac12K[s,s],\qquad
 P=ds+\tfrac12[s,s],\qquad r=\operatorname{pr}P,
 \qquad V=(1+K\operatorname{ad}_s)^{-1}i_2.
\end{equation}
The inverse exists in the \((h)\)-adic topology.  The contraction
identity and Bianchi identity give
\(P=i_2r-K[s,P]\), and hence
\begin{equation}
\label{eq:transfer-mc}
                 P=Vr,\qquad a=s+Vq,\qquad Qh=0,
                 \qquad Qq=-r.
\end{equation}
Here and below the suspension signs in the integrated symplectic form
are fixed compatibly with the last convention.

For completeness, these formulas give the entire transferred
Maurer--Cartan map.  Put \(W=dV+[s,V]\).  The recursion for \(V\) gives
\(KV=0\), \(\operatorname{pr}V=1\), and
\[
 KW=KdV+K[s,V]=V-i_2+K[s,V]=0.
\]
The columns of \(W\) are three-forms in \(C\otimes\mathfrak g\).
They have zero differential by dimension and zero cohomology projection
because \(H^3(C)=0\); the contraction identity therefore gives \(W=0\).
The bracket of two columns of \(Vq\) is a physical four-form and is
zero.  Equations~\eqref{eq:transfer-mc} now imply
\[
                    (d+Q)a+\tfrac12[a,a]=0.
\]
In particular there are no omitted terms quadratic in \(q\).
The induced map of formal moduli problems has the identity on tangent
cohomology, and is a formal equivalence.  One can also see this directly
by induction over nilpotent coefficient ideals: each square-zero
extension has the same cochain deformation and obstruction groups, and
the contraction identifies those groups.

Pull back the strict integrated closed form by this map and denote it
by \(\omega\).  It consists of \(dh\,dq\) terms and
\(q\,dh\,dh\) terms.  Its constant \(dh\,dq\) matrix is the perfect
pairing \(H^1\otimes H^2\to\mathbb C\), so it is formally
nondegenerate.  Explicitly, give coordinate forms total parity equal
to internal degree plus exterior degree, and extend coordinate
differentiation across a physical form \(c\) by
\(\delta(c\otimes u)=(-1)^{|c|}c\otimes\delta u\).
The mixed part is then
\[
 \omega_{\mathrm{mixed}}
   =-\sum_{i,\mu}\left\langle\partial_i s,V_\mu\right\rangle
                                                  dh_i\,dq_\mu.
\]
The contraction \(\iota_Q\) is total-even.  Since
\(\iota_Qdq_\mu=-r_\mu\), contraction removes the \(dq\) and gives
the pairing with \(Vr=P\); its value on the \(q\,dh\,dh\) terms
is zero.  The variation formula for the central
Chern--Simons functional, whose boundary term vanishes, consequently
gives the exact identity
\begin{equation}
\label{eq:transferred-hamiltonian}
                 \iota_Q\omega
                   =\int_Y\operatorname{tr}(\delta s\wedge P)
                   =\delta f.
\end{equation}
The symbol \(\delta\) differentiates only the formal coordinates.
This proves the Hamiltonian identity for the chosen, possibly
noncyclic, contraction.  Reversing the global symplectic sign reverses
both algebraic and Chern--Simons Hamiltonians and has no effect on the
statement of Theorem~\ref{thm:edge-algebraicity}.

\subsection{A closed-form homotopy gives a squared-Jacobian error}

We now have two potentials on equivalent minimal formal models: the
algebraic Darboux potential and the Chern--Simons potential.  The
comparison already identifies their critical schemes and based
holonomies.  We use its additional full closed-form homotopy to place
the difference of potentials in the square of the Jacobian ideal.
This stronger conclusion permits a coordinate correction that fixes
the critical scheme.

Compare the two minimal completed semifree models, the transferred model
and the Darboux model \eqref{eq:algebraic-darboux-chart}, over the formal
completion of \(\mathcal X\).  A map between these models is represented
by a formal power-series map: equivalently one uses homotopy transfer
and maps of minimal Lie models, or lifts maps of completed semifree
algebras order by order.  An equivalence has invertible linear part on
generators, since the cotangent differentials at the point are zero.
It therefore has a formal inverse.  Degree considerations make its
shape particularly simple:
\begin{equation}
\label{eq:minimal-model-isomorphism}
                       x=x(h),\qquad p=B(h)q,
\end{equation}
where the Jacobian of \(x(h)\) and the matrix \(B(h)\) are invertible.
Let \(f'=F(x(h))\), and let \(\omega'\) be the pullback of the
canonical Darboux form.  Then
\[
 \iota_Q\omega'=\delta f',\qquad
 (Qq_1,\ldots,Qq_N)=\Jac(f')=\Jac(f)=J\subset(h)^2.
\]
The first equality is naturality of the Darboux Hamiltonian identity;
the equality of ideals follows from the two invertible matrices in
\eqref{eq:minimal-model-isomorphism} and from the edge curvature
calculation.  The comparison is over \(\mathcal X\), so it also
identifies all the marked classical holonomy functions.

Both strict forms represent the same shifted closed-form class.
We record explicitly the consequence for the Hamiltonians; equality
of the ordinary leading two-form alone would not suffice.

\begin{lemma}[Completed Cartan calculation]
\label{lem:cartan-jacobian-square}
Suppose that \(Qh=0\), \(Qq\in(h)^2\), and that strict closed
\((-1)\)-shifted two-forms \(\omega,\omega'\) on \(R[q]\) are
homotopic as closed forms.  If
\(\iota_Q\omega=\delta f\) and
\(\iota_Q\omega'=\delta f'\), then, after normalizing constants,
\[
                       f-f'\in(Qq_1,\ldots,Qq_N)^2.
\]
\end{lemma}

\begin{proof}
In this proof put \(J=(Qq_1,\ldots,Qq_N)\subset R\).
Use total signs for which \(I=\iota_Q\) is even and
\[
                 [I,\delta]=\mathcal L_Q,
                 \qquad [I,\mathcal L_Q]=0.
\]
The closed-form homotopy has components \(\psi_j\), of exterior
degree \(j\) and internal degree \(-j\), for \(j\geq2\).
Changing the component signs if necessary, its equations are
\[
 \omega-\omega'=\mathcal L_Q\psi_2,
                  \qquad \delta\psi_j=\mathcal L_Q\psi_{j+1}.
\]
Set
\[
             S_j=I^j\psi_j,\qquad
             T_j=I^{j-1}\mathcal L_Q\psi_j.
\]
Commuting \(\delta\) past the \(j\) contractions gives
\(\delta S_j=T_{j+1}-jT_j\), while
\(T_2=\delta(f-f')\).  Thus, for every \(L\geq2\),
\begin{equation}
\label{eq:cartan-factorial-tail}
 \delta(f-f')=-\delta\sum_{j=2}^{L}\frac{S_j}{j!}
                                      +\frac{T_{L+1}}{L!}.
\end{equation}
Every contraction replaces a \(dq_i\) by \(Qq_i\) and kills a
\(dh_i\).  A monomial of \(\psi_j\) containing \(a\) factors \(q\),
\(b\) factors \(dq\), and \(c\) factors \(dh\) satisfies
\(a+b=j=b+c\), by internal and exterior degree.  Hence \(a=c\).
Survival after \(j\) contractions requires \(c=0\), so also \(a=0\)
and \(b=j\).  Consequently \(S_j\) is an ordinary function in \(J^j\).
For the remainder, the surviving monomials contain either \(L+1\)
factors from \(J\), or \(L\) such factors and one differentiated
factor \(\delta(Qq_i)\).  Since
\(Qq_i\in(h)^2\) and \(\delta(Qq_i)\in(h)\,dh\), their coefficient
order is at least \(2L+1\).  Hence the remainder tends to zero.
The sum converges in the complete power-series ring, and
\eqref{eq:cartan-factorial-tail} implies
\[
                 f-f'=-\sum_{j\geq2}\frac{S_j}{j!}+\text{constant}.
\]
The ideal \(J^2\) is adically closed.  The normalized constant is
zero, proving the claim without an isolated-critical-point hypothesis.
\end{proof}

\begin{lemma}[Removing the error while fixing the critical scheme]
\label{lem:jacobian-square-flow}
Let \(f\in(h)^3\) and \(e\in\Jac(f)^2\).
There is a formal automorphism \(\Phi\) of \(R\), equal to the
identity modulo \(J=\Jac(f)\), such that
\[
                           (f+e)\circ\Phi=f.
\]
\end{lemma}

\begin{proof}
Write \(f_i=\partial_i f\) and
\(e=\sum_i v_i f_i\), with \(v_i=\sum_j a_{ij}f_j\in J\).
Since the Hessian of \(f\) has entries in \((h)\), differentiating
this expression gives a matrix \(D\) with entries in \((h)\) such
that
\[
                   \nabla e=D\nabla f.
\]
Explicitly,
\(D_{kj}=\partial_kv_j+\sum_i a_{ij}\partial_k f_i\), which has
the stated order.  Put \(f_t=f+te\).  The vector field
\[
                  X_t=-(1+tD)^{-\mathsf T}v
\]
has entries in \(J\subset(h)^2\) and satisfies
\(df_t(X_t)=-e\).  Its formal flow \(\Phi_t\), with
\(\Phi_0=\id\), exists coefficient by coefficient: each fixed
\((h)\)-adic order involves finitely many powers of \(t\), and
integration of those polynomials is defined over \(\mathbb C\).
It obeys
\[
                 \frac{d}{dt}\Phi_t^*f_t
                    =\Phi_t^*(e+df_t(X_t))=0.
\]
Every \(X_t\) preserves \(J\), since its coefficients belong to
\(J\); its induced derivation on \(R/J\) is zero.  Therefore the
flow preserves \(J\) and is the identity on \(R/J\).
Taking \(\Phi=\Phi_1\) proves both assertions.
\end{proof}

\begin{proof}[Proof of Theorem~\ref{thm:edge-algebraicity}]
Lemma~\ref{lem:cartan-jacobian-square} gives \(f'-f\in J^2\).
Apply Lemma~\ref{lem:jacobian-square-flow} with \(e=f'-f\).
Then
\[
                    f= f'\circ\Phi=F\circ x\circ\Phi,
\]
so \(\varphi=x\circ\Phi\) is the required coordinate change.
The original comparison \(x\) identifies the critical scheme with
the algebraic critical chart and carries its regular holonomy
markings to the specified based holonomies.  For each such entry
\(H\), the additional correction obeys
\[
                      \Phi^*x^*H=x^*H\quad\text{in }R/J,
\]
because \(\Phi\) is the identity on that quotient.  The algebraic
lifts constructed after \eqref{eq:algebraic-darboux-chart} consequently
remain the required lifts after this correction.  This proves the
marking assertion scheme-theoretically.
\end{proof}

\subsection{Polynomial replacement of the active logarithm}

Make the fixed coordinate change of the theorem separately on each
edge, and again write \(f\) for its now algebraic central potential.
For an active edge choose an algebraic matrix lift \(G_\beta\) of
the based \(\beta\)-holonomy, with \(G_\beta(0)=1\).  It suffices to
lift its entries; all identities with the actual holonomy below are
identities in \(R/\Jac(f)\).  Set
\[
 I_\beta=((G_\beta-1)_{ij})\subset R,
 \qquad
 L_d(G_\beta)=\sum_{k=1}^{d}\frac{(-1)^{k+1}}{k}(G_\beta-1)^k.
\]
With the connection convention \(d+a\) of the edge construction,
\(G_\beta=\exp(-v)\) on the central critical scheme.  Its active
coefficient satisfies
\[
             b\equiv-\epsilon\operatorname{tr}
                           (T_0\log G_\beta)\pmod{\Jac(f)}.
\]
Define the algebraic germ
\(b_{\mathrm{alg}}=-\epsilon\operatorname{tr}(T_0 L_d(G_\beta))\).
For every \(d\geq2\), the \((h)\)-adically convergent logarithm series gives
\begin{equation}
\label{eq:polynomial-log-error}
             b-b_{\mathrm{alg}}\in\Jac(f)+I_\beta^{d+1}
                                      \subset\Jac(f)+I_\beta^3.
\end{equation}
If \(b-b_{\mathrm{alg}}=\sum_i c_i\partial_i f+r\), with
\(r\in I_\beta^{d+1}\), then the square-zero coordinate change
\(x_i\mapsto x_i-\tau c_i\) transforms
\(f+\tau b\) into
\[
                         f+\tau(b_{\mathrm{alg}}+r),
                         \qquad \tau^2=0.
\]
Thus the Jacobian part is removed exactly.  The remaining error has
at least cubic order in the multiplicative holonomy entries; it is
this explicit error, rather than an algebraicity assertion about the
formal logarithm, that will be used in the vertex matching.
